% 00-frontmatter.tex — Foreword and Introduction (non-normative)

% =========================================================================
% FOREWORD (non-normative)
% =========================================================================
\chapter*{Foreword}
\addcontentsline{toc}{chapter}{Foreword}

This document specifies the Lain programming language, Version~1.0.

Lain is a statically-typed, compiled programming language for systems
programming.  It guarantees memory safety, type safety, and resource safety
entirely at compile time with zero runtime overhead.

This specification is structured analogously to ISO/IEC~9899 (the C
standard).  Clauses~1 through~4 are administrative; Clauses~5 through~20
are normative.  Annexes~A and~B are normative; Annexes~C and~D are
informative.

Text marked \textbf{Constraint} describes a condition whose violation
requires an implementation to issue a diagnostic message.  Text marked
\idbehav{} denotes behavior that is valid but whose exact result is
determined by the implementation and shall be documented by it.  Text
marked \undefbehav{} denotes a situation for which this specification
imposes no requirements.

\begin{note}
  Undefined behavior in Lain can arise only inside \lain{unsafe} blocks or
  through direct C interoperability.  All code in the safe fragment of Lain
  has defined behavior.
\end{note}

% =========================================================================
% INTRODUCTION (non-normative)
% =========================================================================
\chapter*{Introduction}
\addcontentsline{toc}{chapter}{Introduction}

\section*{Design philosophy}

Lain is built on five foundational pillars:

\begin{description}[leftmargin=2em]
  \item[P1 --- Assembly-Speed Performance]
    Every language feature compiles to C code that is equivalent in
    performance to hand-written C.  There are no hidden allocations, no
    garbage collector, and no reference-counting overhead.

  \item[P2 --- Zero-Cost Memory Safety]
    Memory safety is enforced entirely at compile time through a linear type
    system, a borrow checker, and value range analysis.  A conforming
    implementation that rejects all constraint violations guarantees, for
    safe programs, the absence of use-after-free, double-free, null
    dereference, and buffer overflow at runtime.

  \item[P3 --- Static Verification]
    The verification properties required by this specification are decidable
    and polynomial in program size.  No SMT solver or unbounded fixpoint
    iteration is required.

  \item[P4 --- Determinism]
    Pure functions (\lain{func}) are referentially transparent and
    guaranteed to terminate.  The same inputs always produce the same
    outputs.

  \item[P5 --- Syntactic Simplicity]
    Lain has no implicit copies, no hidden conversions (beyond the defined
    integer-widening rules), and no magic methods.  Ownership is always
    explicit.
\end{description}

\section*{Relationship to C}

Lain compiles to C99.  The C code generated by a conforming Lain
implementation is the \emph{translation output}; this specification does not
prescribe the form of that output beyond the observable behavior of the
resulting program.

\section*{How to read this specification}

Each normative clause is structured as follows: a prose description of the
syntax (with a reference to the formal grammar in Annex~A), followed by the
semantics, followed by \textbf{Constraint} blocks listing every condition
that requires a diagnostic message.  Non-normative \textbf{Note} and
\textbf{Example} environments appear throughout to aid understanding.
